% REPLACES: online_appendix.tex lines 749-770 ("For Proposition A.6, let ..." to "... imposed separately when differentiating a seller objective.")
% The logistic part of A.6 (lines 709-747) and the nonemptiness and lemma parts
% (lines 772-831) are outside this fragment; see M2_20 for their cost lines.
% No \NUM placeholders. Component notation H_L, H_H is kept from the current text:
% H_L is the component with weight \rho, now supported from zero.
For Proposition A.6, let \(H_C=\rho H_L+(1-\rho)H_H\), with atomless
component laws supported within
\([0,\varepsilon_C]\) and
\([c_H-\varepsilon_C,c_H+\varepsilon_C]\). Retain (A3) and impose
\(\varepsilon_C<B_{r_1}(m)\), and

\[
B_{r_0}(1/2)<c_H-\varepsilon_C<c_H+\varepsilon_C<B_{r_1}(M).
\tag{OA.26}\label{eq:oa-oa-26}
\]

Use either the Laplace or logistic noise law. The entry rule is
\(e(\mu)=H_C(B_r(\mu))\). It is continuous and nondecreasing, and
\(e\ge\rho\) because the whole first support lies below gross profit at the
lowest feasible belief in both economies. Equation
\eqref{eq:oa-oa-9} remains strictly positive even if the component
densities are discontinuous or vanish on subintervals. In the weak
economy, \(e(1/2)=\rho\). In the strong economy, a posterior
sufficiently close to \(M\) satisfies \(B_r(\mu)>c_H+\varepsilon_C\); an
event with positive probability then gives \(e=1\). All global trading
comparisons are unchanged. This proof requires an atomless distribution,
not a differentiable density. Differentiability is imposed separately
when differentiating a seller objective.
